\chapter{THIẾT KẾ CÁC LỚP CHI TIẾT}

\section{Bảng mô tả các lớp}

\begin{longtable}{|p{3cm}|p{3cm}|p{6.5cm}|}
\hline
\textbf{Lớp} & \textbf{Thuộc tính chính} & \textbf{Mô tả} \\
\hline
\endfirsthead
\hline
\textbf{Lớp} & \textbf{Thuộc tính chính} & \textbf{Mô tả} \\
\hline
\endhead
NguoiDung & MaNguoiDung, TenDangNhap, MatKhau, VaiTro, TrangThai & Lưu thông tin tài khoản và phân quyền. \\
\hline
HangHoa & MaHangHoa, TenHangHoa, MaNhomHang, MaDonViTinh, GiaNhap, GiaXuat, MoTa, TrangThai & Lưu thông tin hàng hóa. \\
\hline
NhomHangHoa & MaNhomHang, TenNhomHang, MoTa & Phân loại hàng hóa. \\
\hline
DonViTinh & MaDonViTinh, TenDonViTinh & Đơn vị đo lường hàng hóa. \\
\hline
NhaCungCap & MaNhaCungCap, TenNhaCungCap, DiaChi, SoDienThoai, Email & Thông tin nhà cung cấp. \\
\hline
KhachHang & MaKhachHang, TenKhachHang, DiaChi, SoDienThoai, Email & Thông tin khách hàng. \\
\hline
Kho & MaKho, TenKho, DiaChi, NguoiPhuTrach & Thông tin kho hàng. \\
\hline
PhieuNhapKho & MaPhieuNhap, MaNhaCungCap, MaKho, MaNguoiDung, NgayNhap, TongTien, TrangThai & Thông tin phiếu nhập kho. \\
\hline
ChiTietPhieuNhap & MaChiTiet, MaPhieuNhap, MaHangHoa, SoLuong, DonGia, ThanhTien & Chi tiết hàng hóa trong phiếu nhập. \\
\hline
PhieuXuatKho & MaPhieuXuat, MaKhachHang, MaKho, MaNguoiDung, NgayXuat, TongTien, TrangThai & Thông tin phiếu xuất kho. \\
\hline
ChiTietPhieuXuat & MaChiTiet, MaPhieuXuat, MaHangHoa, SoLuong, DonGia, ThanhTien & Chi tiết hàng hóa trong phiếu xuất. \\
\hline
PhieuKiemKe & MaPhieuKiemKe, MaKho, MaNguoiDung, NgayKiemKe, TrangThai & Thông tin phiếu kiểm kê. \\
\hline
ChiTietKiemKe & MaChiTiet, MaPhieuKiemKe, MaHangHoa, SoLuongThucTe, SoLuongSoSach, ChenhLech & Chi tiết kiểm kê. \\
\hline
TonKho & MaKho, MaHangHoa, SoLuongTon, NgayCapNhat & Số lượng tồn kho của từng hàng tại từng kho. \\
\hline
LichSuThaoTac & MaLog, MaNguoiDung, HanhDong, ThoiGian, DoiTuong & Lưu lịch sử thao tác của người dùng. \\
\hline
\end{longtable}

\section{Bảng mô tả quan hệ giữa các lớp}

\begin{longtable}{|p{4cm}|p{4cm}|p{4.5cm}|}
\hline
\textbf{Lớp nguồn} & \textbf{Lớp đích} & \textbf{Loại quan hệ} \\
\hline
\endfirsthead
\hline
\textbf{Lớp nguồn} & \textbf{Lớp đích} & \textbf{Loại quan hệ} \\
\hline
\endhead
NguoiDung -- PhieuNhapKho & 1..n & Một người dùng lập nhiều phiếu nhập. \\
\hline
NguoiDung -- PhieuXuatKho & 1..n & Một người dùng lập nhiều phiếu xuất. \\
\hline
NhaCungCap -- PhieuNhapKho & 1..n & Một nhà cung cấp có nhiều phiếu nhập. \\
\hline
KhachHang -- PhieuXuatKho & 1..n & Một khách hàng có nhiều phiếu xuất. \\
\hline
Kho -- PhieuNhapKho & 1..n & Một kho nhận nhiều phiếu nhập. \\
\hline
Kho -- PhieuXuatKho & 1..n & Một kho xuất nhiều phiếu xuất. \\
\hline
Kho -- TonKho & 1..n & Một kho có nhiều bản ghi tồn kho. \\
\hline
HangHoa -- TonKho & 1..n & Một hàng hóa có tồn kho tại nhiều kho. \\
\hline
HangHoa -- ChiTietPhieuNhap & 1..n & Một hàng hóa xuất hiện trong nhiều chi tiết nhập. \\
\hline
HangHoa -- ChiTietPhieuXuat & 1..n & Một hàng hóa xuất hiện trong nhiều chi tiết xuất. \\
\hline
PhieuNhapKho -- ChiTietPhieuNhap & 1..n & Một phiếu nhập có nhiều dòng chi tiết. \\
\hline
PhieuXuatKho -- ChiTietPhieuXuat & 1..n & Một phiếu xuất có nhiều dòng chi tiết. \\
\hline
NhomHangHoa -- HangHoa & 1..n & Một nhóm hàng có nhiều hàng hóa. \\
\hline
DonViTinh -- HangHoa & 1..n & Một đơn vị tính được dùng cho nhiều hàng hóa. \\
\hline
\end{longtable}

\section{Kết luận chương}

Chương 6 đã mô tả chi tiết các lớp trong hệ thống và các quan hệ giữa chúng. Các thiết kế này là nền tảng để xây dựng mô hình cơ sở dữ liệu và triển khai phần mềm trong Chương 7.
